% ===========================================================================
%  90 — Backlog: o que está pela metade ou por fazer
% ===========================================================================
\chapter{Backlog: o que falta}
\label{cap:backlog}

Este capítulo junta, num lugar só, tudo o que as histórias marcaram como
\parcial{} ou \afazer{}, mais o que o levantamento encontrou fora das telas.
Está organizado por \textbf{horizonte}, não por épico. A coluna \enquote{Quem}
diz de onde vem o trabalho: nem tudo é desenvolvimento.

\begin{tabularx}{\linewidth}{@{}L{22mm}X@{}}
\toprule
\headrow \thd{Quem} & \thd{Significado}\\
\midrule
Dev & Trabalho do time de desenvolvimento.\\
\zebra Conteúdo & Material pedagógico ou de mídia: áudio, foto, revisão de
  guia, teste em sala.\\
Jurídico & Contrato, termo de tratamento de dados, documentos legais.\\
\zebra Produto & Decisão do PO antes de qualquer desenvolvimento (ver o
  Capítulo~\ref{cap:decisoes}).\\
Operação & Configuração ou publicação em cada rede.\\
\zebra Externo & Serviço contratado fora da Unooria, como a revisão de
  segurança.\\
\bottomrule
\end{tabularx}

\section{Bloqueia a entrada em produção (novembro de 2026)}

\begin{tabularx}{\linewidth}{@{}L{15mm}X L{20mm}@{}}
\toprule
\headrow \thd{Ref.} & \thd{O que falta} & \thd{Quem}\\
\midrule
\endhead
\id{H12.4.1} & Teste de carga: 2 a 4 mil pessoas entrando e abrindo missão em
  10 minutos, com turmas inteiras usando IA ao mesmo tempo. Nunca foi feito. &
  Dev\\
\zebra \id{H12.5.1} & Monitoramento e alertas em homologação e produção (hoje
  só em desenvolvimento), com teste do canal de e-mail e alerta de custo. &
  Dev\\
\id{H12.3.1} & Prazos de guarda e eliminação do dado de criança, e a rotina
  que os cumpre. & Produto, Jurídico, Dev\\
\zebra --- & Revisão de segurança externa e teste de intrusão. & Externo\\
\id{H09.3.2} & Contrato e termo de tratamento de dados de cada rede declarando
  que a IA processa fora do Brasil. & Jurídico\\
\zebra \id{H12.2.1} & Ferramenta de gravação de sessão de terceiros nas telas:
  declarar no termo de cada rede ou retirar das telas de menores (\id{D13}). &
  Produto, Jurídico\\
\id{H01.5.1} & Login único da rede (WSO2) validado em homologação. & Dev,
  Operação\\
\zebra \id{H11.3.1} & Razão social e CNPJ nos documentos legais. & Jurídico\\
\bottomrule
\end{tabularx}

\section{Esperado na primeira rede}

\begin{tabularx}{\linewidth}{@{}L{15mm}X L{20mm}@{}}
\toprule
\headrow \thd{Ref.} & \thd{O que falta} & \thd{Quem}\\
\midrule
\endhead
\id{H07.1.4} & Texto livre só para o professor também no Rastro do Boato e no
  Manifesto da IA (\id{D03}). & Dev\\
\zebra \id{H09.4.2} & IA do Instanooria e da Vila do Lume dentro do limite
  diário. & Dev\\
\id{H12.6.1} & Acessibilidade verificada (contraste, foco, teclado, leitor de
  tela) --- costuma ser exigência de edital. & Dev\\
\zebra \id{H01.4.1} \id{H01.4.2} & Trocar a senha por conta própria e recuperar
  a senha esquecida (\id{D05}). & Dev, Produto\\
\id{H03.4.2} & Escolher a escola numa lista, em vez de digitar o código. & Dev\\
\zebra \id{H07.2.1} & Nome da missão (e não o código) no relatório da
  diretoria. & Dev\\
\id{H07.3.1} & Análise de produto no menu, com botão de exportar e seleção de
  escola. & Dev\\
\zebra \id{H11.1.2} & Botão \enquote{Entrar} e links para os documentos legais
  no site; link da Política de Privacidade no cadastro (\id{H11.2.1}). & Dev\\
\id{H01.3.1} & Trocar de escola também no celular. & Dev\\
\zebra \id{H03.3.1} & Filtro Papel completo em Usuários. & Dev\\
\id{H01.2.2} & Retirar o item \enquote{Iniciativas}, que não faz nada, dos
  menus laterais. & Dev\\
\zebra \id{H06.1.1} & Fotos reais e locuções do O Que a Foto Conta. & Conteúdo\\
\id{H06.7.4} \id{H06.4.4} & Áudios do Manifesto da IA e do Rastro do Boato. &
  Conteúdo\\
\zebra \id{H06.8.4} \id{H04.4.1} & Revisão pedagógica e teste em sala do
  guia do unooGPT; autor e revisor do gabarito do Rastro. &
  Conteúdo\\
\id{H10.5.2} & Eliminar um cadastro de prospecção pela tela. & Dev\\
\zebra \id{H06.10.2} \id{H06.12.1} \id{H06.12.4} \id{H06.12.5} & A Seleção (7º ano,
  Unidade 1, \id{F06.12}), já construída e percorrida pelo teste: revisar as
  imagens geradas, conferir com a IA de verdade os níveis que ela dá à legenda
  e revisar os textos escritos pela equipe. & Conteúdo, Dev, Produto\\
\id{H06.10.5} \id{H06.12.6} & Revisar a voz gerada do podcast da A
  Seleção; se não passar, gravar o episódio em estúdio, com o mesmo roteiro. &
  Conteúdo\\
\zebra \id{H06.10.3} \id{H06.11.1} \id{H06.11.4} \id{H06.11.5} & Quem mandou isso?
  versão nova (5º ano, Unidade 1, \id{F06.11}): revisar os textos escritos pela
  equipe e decidir o equilíbrio do pacote pronto (\id{D14}). & Dev, Conteúdo,
  Produto\\
\id{H06.10.4} & Jornal da Escola no lugar do UnooTube (4º ano, Unidade
  1), para o 4º ano voltar a ter missão. & Dev, Conteúdo\\
\zebra \id{H11.3.1} & Revisar a frase dos Termos de Uso que ainda cita o UnooTube,
  arquivado, como aplicativo imitado. & Produto\\
\id{H14.2.3} \id{H14.3.1} \id{H14.3.2} \id{H14.3.3} \id{H14.4.2}
  \id{H14.4.3} \id{H14.6.1} & Teste de ponta a ponta da evidência pedagógica
  que passe pelo que o de hoje não percorre: as equipes, a aula projetada
  (\enquote{a revisar} e \enquote{Confirmar para a turma}), o ajuste de vários
  estudantes e o histórico, o \enquote{Não participou}, a ordem da lista e um
  perfil com conteúdo. & Dev\\
\zebra \id{F13.4} \id{H13.4.1} \id{H13.4.2} \id{H13.4.3} \id{H13.4.4} & Equipes do
  mesmo lado do feed no Instanooria: o sorteio agrupa pelo lado, quem entra
  recebe o lado da equipe, a redistribuição fica desligada e o monitor mostra
  as equipes. É a próxima frente de \id{E13}. & Dev\\
\bottomrule
\end{tabularx}

\section{Evolução}

\begin{tabularx}{\linewidth}{@{}L{15mm}X L{20mm}@{}}
\toprule
\headrow \thd{Ref.} & \thd{O que falta} & \thd{Quem}\\
\midrule
\endhead
\id{H04.2.3} & Liberar a missão com data de abertura e fechamento (\id{D12}). &
  Dev\\
\zebra \id{H10.1.3} & Desligar recurso por turma ou até uma data. & Dev\\
\id{H10.3.3} & Aviso de versão fixada há muito tempo. & Dev\\
\zebra \id{H02.4.2} & Virada do ano assistida (a primeira real é em fevereiro de
  2027). & Dev\\
\id{H09.4.3} & Aviso quando o consumo de IA de uma rede chegar a 80\%. & Dev\\
\zebra \id{H07.3.2} & Números do dia para a rede. & Dev\\
\id{H05.2.2} & Guardar os registros de uso sem conexão. & Dev\\
\zebra \id{F06.10} \id{H06.10.1} & As 41 unidades da trilha ainda em produção,
  cada uma entrando pelo mesmo caminho, já completa (\id{H06.10.1}). &
  Conteúdo, Dev\\
\id{H08.3.4} & Visão agregada e anônima do bem-estar para a rede --- só se
  decidida (\id{D11}). & Produto, Jurídico\\
\zebra \id{H14.1.1} & A régua pedagógica nas outras missões. A primeira ---
  Quem mandou isso?, versão nova --- já entrou e tornou visível ao professor o
  que \id{E14} faz, e A Seleção veio em seguida; o Instanooria é a próxima
  candidata. & Conteúdo, Dev\\
\id{F14.1} & O unooGPT refeito com atividades, para entrar na
  evidência pedagógica. & Dev, Conteúdo\\
\zebra \id{H14.4.4} & A tela em que a coordenação consulta a evidência (a
  leitura já é permitida). & Dev\\
\id{H14.6.2} & O painel que junta as habilidades de todas as missões, com a
  regra definida pela equipe pedagógica. & Produto, Conteúdo, Dev\\
\bottomrule
\end{tabularx}
